%% Springer Nature LaTeX template (sn-jnl.cls v3.1, December 2024)
%% Math and Physical Sciences, Numbered reference style.
%% Companion paper "On the non-existence of non-trivial Collatz cycles..."
%% by Eric Merle, prepared for submission to Journal of Automated Reasoning.

\documentclass[pdflatex,sn-mathphys-num]{sn-jnl}

%%%% Standard packages
\usepackage{graphicx}
\usepackage{multirow}
\usepackage{amsmath,amssymb,amsfonts}
\usepackage{amsthm}
\usepackage{mathtools}
\usepackage{mathrsfs}
\usepackage[title]{appendix}
\usepackage{xcolor}
\usepackage{textcomp}
\usepackage{manyfoot}
\usepackage{booktabs}
\usepackage{array}
\usepackage{longtable}
\usepackage{calc} % for tabularnewline width arithmetic Pandoc emits
\usepackage{listings}
\usepackage{tikz}
\usetikzlibrary{positioning, arrows.meta, shapes.geometric}
\usepackage{hyperref}

%%%% Pandoc-emitted commands the body relies on
\providecommand{\tightlist}{\setlength{\itemsep}{0pt}\setlength{\parskip}{0pt}}
\providecommand{\citeproc}[2]{[#1]} % no-op when natbib is in use

% Pandoc emits "\def\LTcaptype{none}" for caption-less longtables; the
% caption package would then look up a counter named "none" that does
% not exist. Define a stub counter so caption-less longtables compile.
\newcounter{none}

%%%% Theorem-like environments — numbered SECTION-WISE so that the
%%%% in-text references "Theorem 3.1", "Theorem 7.1", "Lemma 5.1" used
%%%% throughout the body match the rendered numbering. This is the
%%%% convention inherited from the Pandoc/amsart draft and from the
%%%% Markdown sources.
\theoremstyle{thmstyleone}
\newtheorem{theorem}{Theorem}[section]
\newtheorem{lemma}[theorem]{Lemma}
\newtheorem{corollary}[theorem]{Corollary}
\newtheorem{proposition}[theorem]{Proposition}

\theoremstyle{thmstyletwo}
\newtheorem{example}{Example}[section]
\newtheorem{remark}{Remark}[section]
\newtheorem*{note}{Note}

\theoremstyle{thmstylethree}
\newtheorem{definition}{Definition}[section]
\newtheorem{conjecture}[theorem]{Conjecture}
\newtheorem{reformulation}[theorem]{Reformulation}

%%%% Unicode glyph fallback (the Markdown source uses Unicode arrows/relations).
\usepackage{newunicodechar}
% Arrows
\newunicodechar{→}{\ensuremath{\rightarrow}}
\newunicodechar{←}{\ensuremath{\leftarrow}}
\newunicodechar{↔}{\ensuremath{\leftrightarrow}}
\newunicodechar{⇒}{\ensuremath{\Rightarrow}}
\newunicodechar{⇐}{\ensuremath{\Leftarrow}}
\newunicodechar{⇔}{\ensuremath{\Leftrightarrow}}
% Order / equality relations
\newunicodechar{≤}{\ensuremath{\leq}}
\newunicodechar{≥}{\ensuremath{\geq}}
\newunicodechar{≠}{\ensuremath{\neq}}
\newunicodechar{≈}{\ensuremath{\approx}}
\newunicodechar{≡}{\ensuremath{\equiv}}
\newunicodechar{∼}{\ensuremath{\sim}}
% Set / logic
\newunicodechar{∈}{\ensuremath{\in}}
\newunicodechar{∉}{\ensuremath{\notin}}
\newunicodechar{∪}{\ensuremath{\cup}}
\newunicodechar{∩}{\ensuremath{\cap}}
\newunicodechar{∧}{\ensuremath{\wedge}}
\newunicodechar{∨}{\ensuremath{\vee}}
\newunicodechar{¬}{\ensuremath{\neg}}
\newunicodechar{∀}{\ensuremath{\forall}}
\newunicodechar{∃}{\ensuremath{\exists}}
\newunicodechar{∅}{\ensuremath{\emptyset}}
\newunicodechar{ℕ}{\ensuremath{\mathbb{N}}}
\newunicodechar{ℤ}{\ensuremath{\mathbb{Z}}}
\newunicodechar{ℚ}{\ensuremath{\mathbb{Q}}}
\newunicodechar{ℝ}{\ensuremath{\mathbb{R}}}
\newunicodechar{ℂ}{\ensuremath{\mathbb{C}}}
% Greek lowercase used in body
\newunicodechar{α}{\ensuremath{\alpha}}
\newunicodechar{β}{\ensuremath{\beta}}
\newunicodechar{γ}{\ensuremath{\gamma}}
\newunicodechar{δ}{\ensuremath{\delta}}
\newunicodechar{ε}{\ensuremath{\varepsilon}}
\newunicodechar{ζ}{\ensuremath{\zeta}}
\newunicodechar{η}{\ensuremath{\eta}}
\newunicodechar{θ}{\ensuremath{\theta}}
\newunicodechar{κ}{\ensuremath{\kappa}}
\newunicodechar{λ}{\ensuremath{\lambda}}
\newunicodechar{μ}{\ensuremath{\mu}}
\newunicodechar{ν}{\ensuremath{\nu}}
\newunicodechar{π}{\ensuremath{\pi}}
\newunicodechar{ρ}{\ensuremath{\rho}}
\newunicodechar{σ}{\ensuremath{\sigma}}
\newunicodechar{τ}{\ensuremath{\tau}}
\newunicodechar{φ}{\ensuremath{\varphi}}
\newunicodechar{χ}{\ensuremath{\chi}}
\newunicodechar{ψ}{\ensuremath{\psi}}
\newunicodechar{ω}{\ensuremath{\omega}}
% Greek uppercase
\newunicodechar{Δ}{\ensuremath{\Delta}}
\newunicodechar{Λ}{\ensuremath{\Lambda}}
\newunicodechar{Π}{\ensuremath{\Pi}}
\newunicodechar{Σ}{\ensuremath{\Sigma}}
\newunicodechar{Φ}{\ensuremath{\Phi}}
\newunicodechar{Ψ}{\ensuremath{\Psi}}
\newunicodechar{Ω}{\ensuremath{\Omega}}
% Misc
\newunicodechar{∞}{\ensuremath{\infty}}
\newunicodechar{·}{\ensuremath{\cdot}}
\newunicodechar{×}{\ensuremath{\times}}
\newunicodechar{÷}{\ensuremath{\div}}
\newunicodechar{±}{\ensuremath{\pm}}
\newunicodechar{∓}{\ensuremath{\mp}}
\newunicodechar{√}{\ensuremath{\sqrt{}}}
\newunicodechar{²}{\ensuremath{^{2}}}
\newunicodechar{³}{\ensuremath{^{3}}}
\newunicodechar{⁷}{\ensuremath{^{7}}}
\newunicodechar{⁰}{\ensuremath{^{0}}}
\newunicodechar{¹}{\ensuremath{^{1}}}
\newunicodechar{₀}{\ensuremath{_{0}}}
\newunicodechar{₁}{\ensuremath{_{1}}}
\newunicodechar{₂}{\ensuremath{_{2}}}
\newunicodechar{₃}{\ensuremath{_{3}}}
\newunicodechar{ₖ}{\ensuremath{_{k}}}
\newunicodechar{ₛ}{\ensuremath{_{s}}}
\newunicodechar{ₙ}{\ensuremath{_{n}}}
\newunicodechar{₈}{\ensuremath{_{8}}}
\newunicodechar{₉}{\ensuremath{_{9}}}
\newunicodechar{ⁿ}{\ensuremath{^{n}}}
\newunicodechar{ᵢ}{\ensuremath{_{i}}}
\newunicodechar{ᵤ}{\ensuremath{_{u}}}
\newunicodechar{ᵥ}{\ensuremath{_{v}}}
\newunicodechar{ⱼ}{\ensuremath{_{j}}}
\newunicodechar{ₘ}{\ensuremath{_{m}}}
% Calculus operators
\newunicodechar{∫}{\ensuremath{\int}}
\newunicodechar{∂}{\ensuremath{\partial}}
\newunicodechar{∇}{\ensuremath{\nabla}}
\newunicodechar{∑}{\ensuremath{\sum}}
\newunicodechar{∏}{\ensuremath{\prod}}
% Section sign: keep as text
\newunicodechar{§}{\S}
% Quotation marks
\newunicodechar{«}{\guillemotleft}
\newunicodechar{»}{\guillemotright}

\raggedbottom

\begin{document}

\title[Conditional non-existence of Collatz cycles in Lean 4]{On the non-existence of non-trivial Collatz cycles:\\
a conditional formal proof in Lean 4 with documented structural obstructions}

\author*[1]{\fnm{Eric} \sur{Merle}}\email{eric.merle@ac-versailles.fr}
\equalcont{ORCID: \href{https://orcid.org/0009-0008-7940-402X}{0009-0008-7940-402X}}

\affil*[1]{\orgname{Independent researcher},
  \orgaddress{\city{Chartres}, \country{France}}}

\abstract{We address the no-non-trivial-cycle disjunct of the Collatz conjecture. The central result, formalized in Lean~4 with Mathlib v4.27.0, is a conditional non-existence theorem declared parametrically with three structural hypotheses: \texttt{BakerSeparation} (after Baker 1966), \texttt{BarinaVerification} (after the 2025 computational verification that every positive integer below $2^{71}$ reaches~$1$), and \texttt{ProductBoundThreshold}, a project-derived cycle-complexity bound made fully explicit in §4.3. The Lean formalization is machine-verified with kernel-3 axiom profile, no user axioms and no \texttt{sorry}, and is reproducible via \texttt{reproduce.sh} against an explicit \texttt{expected\_axioms.md} baseline.

We complement the theorem with two impossibility lemmas ($\delta 8$, $\delta 8'$) showing that no uniform algebraic refinement of the Product Bound derivation can eliminate the third hypothesis within the Baker + continued-fraction framework; a literature mapping ($\delta 9$) documenting the absence of any peer-reviewed deterministic upper bound on cycle length in the 1977--2026 results we surveyed; and an alternative disjunctive framing ($\delta 7$) that connects the conditional theorem to Hercher's 2023 lower bound. The paper is intended as a stable substrate for future contributions resolving the structural conditionality.}

\keywords{Collatz conjecture, 3x+1 problem, Collatz cycles,
linear forms in logarithms, Baker's theorem, continued fractions,
irrationality measure, formal verification, Lean 4, Mathlib}

\pacs[MSC Classification]{11B83, 11J81, 11J86, 11Y55, 37P99, 68V15}

\maketitle

%%%%%%%%%%%%%%%%%%%%%%%%%%%%%%%%%%%%%%%%%%%%%%%%%%%%%%%%%%%%%%%%%%%%%%
%% Body content (auto-generated from Markdown sources via Pandoc)  %%
%%%%%%%%%%%%%%%%%%%%%%%%%%%%%%%%%%%%%%%%%%%%%%%%%%%%%%%%%%%%%%%%%%%%%%

\input{body.tex}

%%%%%%%%%%%%%%%%%%%%%%%%%%%%%%%%%%%%%%%%%%%%%%%%%%%%%%%%%%%%%%%%%%%%%%
%% References. The .bib file is the same as the one used by the   %%
%% Pandoc draft build (paper/v2/references.bib). \nocite{*} forces %%
%% inclusion of every entry until the body is converted to use     %%
%% explicit \cite{} commands.                                       %%
%%%%%%%%%%%%%%%%%%%%%%%%%%%%%%%%%%%%%%%%%%%%%%%%%%%%%%%%%%%%%%%%%%%%%%
\nocite{*}
\bibliography{references}

\end{document}
